\subsection{Retrieval-Augmented Generation Architecture}

Our retrieval-augmented generation (RAG) pipeline combines
automatic speech recognition (ASR), dense retrieval, and
evidence-conditioned answer generation. Given an audio query,
the ASR component first produces a transcript, which is
encoded into a dense query representation. The retriever
identifies relevant corpus records, and the generator
produces an answer conditioned on the selected evidence.

\paragraph{Exact dense retrieval.}
We use precomputed, L2-normalised corpus embeddings and
an exact dense retriever. For a query embedding
$\mathbf{q}$ and corpus embedding $\mathbf{d}_i$,
cosine similarity is computed as
\begin{equation}
s_i =
\frac{\mathbf{q}^{\top}\mathbf{d}_i}
{\|\mathbf{q}\|_2\|\mathbf{d}_i\|_2}.
\end{equation}
Because corpus embeddings are already normalised and
query embeddings are normalised at retrieval time,
ranking reduces to a matrix--vector multiplication.
The implementation selects the highest-scoring $K$
records using partial sorting and orders the selected
records by descending similarity.

This exact-search baseline avoids approximation errors
introduced by an approximate nearest-neighbour index,
although retrieval requires computing similarities
against the corpus embeddings.

\paragraph{Context construction and generation.}
The context builder retrieves the text corresponding
to the selected record identifiers and concatenates
the documents into a numbered evidence block.
The resulting prompt contains the transcribed question,
retrieved evidence, and an instruction to acknowledge
insufficient evidence rather than invent an answer.

We distinguish the retrieval depth $K$ from the number
of documents included in the generation context, $N$,
where $1 \leq N \leq K$. This separation allows us
to study the quality--cost trade-off of providing
additional evidence while holding the retrieval
candidate set fixed.

\paragraph{Instrumentation.}
The pipeline records ASR, query embedding, retrieval,
generation, and end-to-end latencies. The generator
interface additionally exposes token-usage information
when returned by the underlying model service.
These measurements support analysis of retrieval and
generation costs alongside answer quality.
